\section{De extremo a extremo: del desarrollo de funciones al entrenamiento de capacidades}

La sección anterior explicó cómo BEV unifica el espacio; esta sección examina un giro mayor: el enfoque de extremo a extremo. Lang Xianpeng (郎咸朋) dijo: «antes hacíamos mal la conducción autónoma». Se refería a desarrollar la conducción autónoma como si fuera una función: enumerar escenarios, escribir reglas, delimitar el ODD y corregir los bugs uno por uno con parches. Pero los escenarios de conducción reales son prácticamente imposibles de enumerar por completo, y las variables tampoco son independientes entre sí: corregir un escenario puede romper otro. El sentido del enfoque de extremo a extremo es convertir la conducción de «reglas escritas a mano» en «capacidades que el modelo aprende de los datos».

\begin{figure}[H]
\centering
\includegraphics[width=0.96\textwidth]{e2e-stack.png}
\caption{Pila de conducción autónoma de extremo a extremo: de la entrada de sensores a la salida de trayectoria, las capas intermedias se van comprimiendo en un sistema que aprende. Diagrama conceptual propio, elaborado a partir de la conversación entre 00:28:46 y 00:41:14.}
\end{figure}

\begin{knowledgebox}{Lectura de la figura: de extremo a extremo no significa sin estructura, sino que el modelo asume más correspondencias}
De Sensors, BEV, Planning y Trajectory hasta Control, en el sistema tradicional cada capa tiene una gran cantidad de reglas escritas a mano. El enfoque de extremo a extremo no elimina por completo la ingeniería, sino que entrega las correspondencias clave al modelo y a los datos, para que el sistema aprenda la capacidad de ir de la entrada a la trayectoria.
\end{knowledgebox}

\subsection{Por qué no es viable enumerar todos los escenarios}

Esta subsección desglosa el fracaso de los sistemas basados en reglas. Las combinaciones de variables como el clima, el tráfico, la estructura de la vía, el estado del pavimento, los vehículos cercanos, los peatones, la velocidad del propio vehículo, las obras, los baches y los objetos extraños son enormes. Aun si se definieran a mano decenas de millones de escenarios, surgirían dos problemas: los escenarios no son ortogonales y las reglas se afectan entre sí; y siempre aparecerán eventos de cola larga que una persona evita con sentido común, pero que un sistema de reglas quizá no sepa cómo manejar.

\begin{warningbox}{El dilema de la cola larga en los sistemas de reglas}
La conducción autónoma no es una lógica de botones. Una regla de vuelta a la derecha se comporta de forma completamente distinta con sol, con lluvia, con tráfico congestionado, con peatones, con bicicletas, con obras o con baches en el pavimento. Cuantas más reglas hay, más graves son las interferencias entre ellas y mayor es el costo de mantenimiento del sistema.
\end{warningbox}

\subsection{Un sistema de capacidades: como enseñar a un niño, no aprender por él}

Esta subsección conserva una metáfora oral muy importante de Lang Xianpeng. Con el enfoque de extremo a extremo, el ingeniero es como un observador y un maestro: proporciona datos de alta calidad, recursos de entrenamiento, preguntas de evaluación y retroalimentación, pero no puede aprender en lugar del modelo. El ingeniero no necesariamente puede explicar del todo en qué se convierte el modelo ni cómo se forman sus estrategias internas; pero puede evaluar si progresa mediante case, pruebas de regresión, simulación, pruebas en carretera y retroalimentación de los usuarios.

\begin{importantbox}{Nota de clase: el modelo y los datos son dos caras de la misma moneda}
Construir un sistema de capacidades con un modelo implica necesariamente depender de los datos; solo las empresas que poseen datos de alta calidad, verticales y con retroalimentación real tienen la oportunidad de mejorar de forma continua la conducción autónoma. Los parámetros del modelo, las técnicas algorítmicas y la experiencia del equipo importan, pero sin datos de conducción de alta calidad el techo será muy bajo.
\end{importantbox}

\subsection{Qué hacen los ingenieros después del enfoque de extremo a extremo}

Esta subsección aclara un punto que se presta a malentendidos. El enfoque de extremo a extremo no significa que los ingenieros se retiren, ni que la responsabilidad de seguridad se entregue a una caja negra. Al contrario, el trabajo del ingeniero pasa de «escribir reglas para cada escenario» a «definir los datos, la evaluación, la regresión y los límites de seguridad». Deben encontrar los case de cola larga y juzgar qué escenarios revelan los puntos ciegos del modelo; deben diseñar el etiquetado y los datos de entrenamiento para que el modelo vea muestras de alto valor; deben construir la simulación y la evaluación en lazo cerrado para evitar que un modelo nuevo corrija un escenario y rompa otro; y deben controlar el ritmo de liberación, para asegurar que los usuarios de vehículos de producción en serie no sean sujetos de experimentación sin protección.

\begin{knowledgebox}{Experiencia práctica: el enfoque de extremo a extremo traslada la dificultad de ingeniería a los datos y la evaluación}
En la era de las reglas, la dificultad de ingeniería estaba en las ramas de código y en la enumeración de escenarios; en la era de los modelos, está en la cobertura de datos, el peso de las muestras, la estabilidad del entrenamiento, las métricas fuera de línea, la experiencia en línea y la regresión de seguridad. El trabajo de ingeniería no desaparece: cambia de lugar.
\end{knowledgebox}

\noindent\begin{tabularx}{\textwidth}{p{0.22\textwidth}p{0.34\textwidth}X}
\toprule
Rol de ingeniería & Trabajo principal en la era de las reglas & Trabajo principal en la era de extremo a extremo \\
\midrule
Ingeniero de algoritmos & Escribir detección, planificación, reglas y máquinas de estados & Diseñar la estructura del modelo, los objetivos de entrenamiento, el muestreo de datos y la evaluación. \\
Ingeniero de datos & Recopilar unos pocos case para validar reglas & Construir las cadenas de procesamiento para el retorno de datos de la flota, la limpieza, la anonimización, el etiquetado y la minería. \\
Ingeniero de pruebas & Ejecutar escenarios fijos y regresiones manuales & Mantener la biblioteca de case de cola larga, la simulación, las métricas fuera de línea y el despliegue gradual en línea. \\
Producto/seguridad & Definir los límites de la función y los avisos & Definir la responsabilidad de la toma de control, el ODD, la degradación ante riesgos y la educación del usuario. \\
\bottomrule
\end{tabularx}

\subsection{De extremo a extremo, RL, VLM y modelos del mundo}

Esta subsección se conecta con la ruta posterior. En la entrevista se discuten juntos el enfoque de extremo a extremo, los VLM, los modelos del mundo y el RL. El enfoque de extremo a extremo hace primero que el sistema aprenda el comportamiento de conducción desde la percepción hasta la trayectoria; los VLM incorporan comprensión semántica e interpretación de escenarios complejos; los modelos del mundo predicen el efecto de las acciones sobre el estado futuro del tráfico; y el RL aporta un marco de optimización en lazo cerrado. La conducción autónoma se parece cada vez más a un Agent en el mundo físico: observar, predecir, actuar, recibir retroalimentación y volver a aprender.

\begin{figure}[H]
\centering
\includegraphics[width=0.96\textwidth]{e2e-rl-worldmodel.png}
\caption{De extremo a extremo + VLM + modelo del mundo: de la imitación de trayectorias al aprendizaje por refuerzo y la predicción de estados futuros. Diagrama conceptual propio, elaborado a partir de la conversación entre 01:15:15 y 01:26:40.}
\end{figure}

\begin{knowledgebox}{Lectura de la figura: lo que aporta el modelo del mundo es la «predicción de consecuencias»}
Una política de extremo a extremo puede producir una trayectoria, pero el tráfico complejo exige predecir: si cambio de carril, reduzco la velocidad o cedo el paso, cómo cambiarán los vehículos y las personas alrededor. El modelo del mundo incorpora al sistema las consecuencias de las acciones, y el RL permite que el sistema optimice su política a partir de la retroalimentación.
\end{knowledgebox}

\subsection{Resumen de la sección}

La esencia del enfoque de extremo a extremo no es una palabra de moda, sino el paso de la conducción autónoma del desarrollo de funciones de software al entrenamiento de capacidades. Convierte al ingeniero, de alguien que escribe reglas a mano, en diseñador de los datos, el entrenamiento, la evaluación y los límites de seguridad.
